\sectionkicker{Source-backed technical notes}
\subsection*{Openは「重み公開」から配備・再現性の設計問題へ: Technical Notes}
本欄は記事本文で圧縮した一次資料上の情報を、比較・再検証しやすい形で整理したものである。時系列、確認済みの主張、留意点、一次資料URLを掲載する。
\smallskip
\noindent{\small\textit{Technical Notesでは、一次資料から確認した公開・提供・研究上の事実を記載する。数値・能力評価や提供元・著者による比較表現は、独立再現済みの結果ではなく帰属claimとして扱う。}}\par
\medskip
\subsection*{Theme at a glance}
\addcontentsline{toc}{subsection}{Theme at a glance}
\begin{center}
\footnotesize
\begin{tabularx}{\linewidth}{@{}p{0.20\linewidth}p{0.10\linewidth}p{0.14\linewidth}X@{}}
\toprule
資料 & 位置づけ & 種別 & 時系列 \\
\midrule
Llama 3.1 & 主要資料 & オープンウェイト & 2024-07-23                   \\
Llama 3.2 & 主要資料 & オープンウェイト & 2024-09-25                   \\
Qwen2.5 & 主要資料 & オープンウェイト & 2024-09-19                   \\
DeepSeek-V3 & 主要資料 & オープンウェイト & 2024-12-26                   \\
OLMoE & 主要資料 & オープンウェイト & 2024-09-04                   \\
OLMo 2 & 主要資料 & オープンウェイト & 2024-11-26                   \\
Tülu 3 & 主要資料 & Framework & 2024-11-21                   \\
Jamba 1.5 Mini and Large & 主要資料 & オープンウェイト & 2024-08-22                   \\
Hunyuan-Large & 主要資料 & オープンウェイト & 2024-11-05                   \\
Phi-3.5 model family & 主要資料 & オープンウェイト & 2024-08-22                   \\
Ministral 3B and 8B & 主要資料 & モデル & 2024-10-16                   \\
Mistral Large 2 & 補足資料 & モデル & 2024-07-24                   \\
Qwen2.5-Coder family & 補足資料 & オープンウェイト & 2024-11-12                   \\
Qwen2.5-Turbo 1M context & 補足資料 & モデル更新 & 2024-11-15                   \\
Quantized Llama 3.2 lightweight models & 補足資料 & モデル更新 & 2024-10-24                   \\
\bottomrule
\end{tabularx}
\end{center}
\normalsize

\begin{technicalnote}{Llama 3.1}{主要資料}
\begingroup
% reader-facing Technical Notes break policy
\widowpenalty=10000
\clubpenalty=10000
\displaywidowpenalty=10000
\begin{tabularx}{\linewidth}{@{}>{\bfseries}p{0.22\linewidth}X@{}}
組織 & Meta \\
種別 & オープンウェイト \\
時系列 & 2024-07-23 \\
\end{tabularx}
\smallskip
{\bfseries 一次資料から整理したtechnical points}
\begin{itemize}[leftmargin=1.5em,itemsep=0.35em]
\item \textbf{一次情報で確認できる事実}: Metaの一次資料により、Llama 3.1について2024-07-23にモデル公開を確認した。 対象event近傍の一次資料から 405B parameter scale を確認できる。
\end{itemize}
\begin{minipage}{\linewidth}
% reader-facing Technical Notes generic boundary/source tail group
\begin{samepage}
{\bfseries 一次資料}
\begingroup\sloppy
\Urlmuskip=0mu plus 2mu\relax
\begin{itemize}[leftmargin=1.5em,itemsep=0.25em]
\item {\scriptsize\url{https://ai.meta.com/blog/meta-llama-3-1/}}
\end{itemize}
\endgroup
\end{samepage}
\end{minipage}
\endgroup
\end{technicalnote}
\medskip

\begin{technicalnote}{Llama 3.2}{主要資料}
\begingroup
% reader-facing Technical Notes break policy
\widowpenalty=10000
\clubpenalty=10000
\displaywidowpenalty=10000
\begin{tabularx}{\linewidth}{@{}>{\bfseries}p{0.22\linewidth}X@{}}
組織 & Meta \\
種別 & オープンウェイト \\
時系列 & 2024-09-25 \\
\end{tabularx}
\smallskip
{\bfseries 一次資料から整理したtechnical points}
\begin{itemize}[leftmargin=1.5em,itemsep=0.35em]
\item \textbf{一次情報で確認できる事実}: Metaの一次資料により、Llama 3.2について2024-09-25にモデル公開を確認した。 Metaの受理済み2024-09-25一次資料では、Llama 3.2を11B / 90Bのvision modelsと1B / 3Bのlightweight text-only modelsとして公開し、1B / 3Bは128K contextをサポートすると記載している。同時発表のRetrieval-Augmented Generation (RAG) / toolingのturnkey deploymentはLlama 3.2 model自身の属性ではなく、Llama Stack distributions / APIsのdeployment capabilityとして分離して扱う。
\end{itemize}
\begin{minipage}{\linewidth}
% reader-facing Technical Notes generic boundary/source tail group
\begin{samepage}
{\bfseries 一次資料}
\begingroup\sloppy
\Urlmuskip=0mu plus 2mu\relax
\begin{itemize}[leftmargin=1.5em,itemsep=0.25em]
\item {\scriptsize\url{https://ai.meta.com/blog/llama-3-2-connect-2024-vision-edge-mobile-devices/}}
\end{itemize}
\endgroup
\end{samepage}
\end{minipage}
\endgroup
\end{technicalnote}
\medskip

\begin{technicalnote}{Qwen2.5}{主要資料}
\begingroup
% reader-facing Technical Notes break policy
\widowpenalty=10000
\clubpenalty=10000
\displaywidowpenalty=10000
\begin{tabularx}{\linewidth}{@{}>{\bfseries}p{0.22\linewidth}X@{}}
組織 & Qwen \\
種別 & オープンウェイト \\
時系列 & 2024-09-19 \\
\end{tabularx}
\smallskip
{\bfseries 一次資料から整理したtechnical points}
\begin{itemize}[leftmargin=1.5em,itemsep=0.35em]
\item \textbf{一次情報で確認できる事実}: Qwenの一次資料により、Qwen2.5について2024-09-19にモデル公開を確認した。 Qwen Teamの受理済み2024-09-19一次資料では、Qwen2.5 familyを0.5B / 1.5B / 3B / 7B / 14B / 32B / 72Bのvariantとして公開している。licenseはvariant単位で区別し、同資料は3Bと72BをApache 2.0の例外と明記しているため、本欄では0.5B / 1.5B / 7B / 14B / 32BのみをApache 2.0対象として扱い、3B / 72Bへ同licenseを結合しない。
\end{itemize}
\begin{minipage}{\linewidth}
% reader-facing Technical Notes generic boundary/source tail group
\begin{samepage}
{\bfseries 一次資料}
\begingroup\sloppy
\Urlmuskip=0mu plus 2mu\relax
\begin{itemize}[leftmargin=1.5em,itemsep=0.25em]
\item {\scriptsize\url{https://qwenlm.github.io/blog/qwen2.5/}}
\end{itemize}
\endgroup
\end{samepage}
\end{minipage}
\endgroup
\end{technicalnote}
\medskip

\begin{technicalnote}{DeepSeek-V3}{主要資料}
\begingroup
% reader-facing Technical Notes break policy
\widowpenalty=10000
\clubpenalty=10000
\displaywidowpenalty=10000
\begin{tabularx}{\linewidth}{@{}>{\bfseries}p{0.22\linewidth}X@{}}
組織 & DeepSeek \\
種別 & オープンウェイト \\
時系列 & 2024-12-26 \\
\end{tabularx}
\smallskip
{\bfseries 一次資料から整理したtechnical points}
\begin{itemize}[leftmargin=1.5em,itemsep=0.35em]
\item \textbf{一次情報で確認できる事実}: DeepSeekの一次資料により、DeepSeek-V3について2024-12-26にモデル公開を確認した。 対象event近傍の一次資料から DeepSeek-V3 served through the existing deepseek-chat API alias を確認できる。
\end{itemize}
\begin{minipage}{\linewidth}
% reader-facing Technical Notes generic boundary/source tail group
\begin{samepage}
{\bfseries 一次資料}
\begingroup\sloppy
\Urlmuskip=0mu plus 2mu\relax
\begin{itemize}[leftmargin=1.5em,itemsep=0.25em]
\item {\scriptsize\url{https://api-docs.deepseek.com/updates}}
\end{itemize}
\endgroup
\end{samepage}
\end{minipage}
\endgroup
\end{technicalnote}
\medskip

\begin{technicalnote}{OLMoE}{主要資料}
\begingroup
% reader-facing Technical Notes break policy
\widowpenalty=10000
\clubpenalty=10000
\displaywidowpenalty=10000
\begin{tabularx}{\linewidth}{@{}>{\bfseries}p{0.22\linewidth}X@{}}
組織 & Ai2 \\
種別 & オープンウェイト \\
時系列 & 2024-09-04 \\
\end{tabularx}
\smallskip
{\bfseries 一次資料から整理したtechnical points}
\begin{itemize}[leftmargin=1.5em,itemsep=0.35em]
\item \textbf{一次情報で確認できる事実}: Ai2の一次資料により、OLMoEについて2024-09-04にモデル公開を確認した。 Ai2の受理済み2024-09-04一次資料では、OLMoEをOLMo family初のsparse Mixture-of-Experts modelとして公開し、1B active / 7B total parameters、5T training tokensと記載している。releaseにはdata、code、evaluations、logs、intermediate training checkpointsを含み、pretrained modelについて5000 stepsごとの244 checkpointsなども公開した。training speed、Pareto frontier、benchmark等の性能表現はAi2による報告として扱い、独立再現結果へ昇格させない。
\end{itemize}
\begin{minipage}{\linewidth}
% reader-facing Technical Notes generic boundary/source tail group
\begin{samepage}
{\bfseries 一次資料}
\begingroup\sloppy
\Urlmuskip=0mu plus 2mu\relax
\begin{itemize}[leftmargin=1.5em,itemsep=0.25em]
\item {\scriptsize\url{https://allenai.org/blog/olmoe-an-open-small-and-state-of-the-art-mixture-of-experts-model-c258432d0514}}
\end{itemize}
\endgroup
\end{samepage}
\end{minipage}
\endgroup
\end{technicalnote}
\medskip

\begin{technicalnote}{OLMo 2}{主要資料}
\begingroup
% reader-facing Technical Notes break policy
\widowpenalty=10000
\clubpenalty=10000
\displaywidowpenalty=10000
\begin{tabularx}{\linewidth}{@{}>{\bfseries}p{0.22\linewidth}X@{}}
組織 & Ai2 \\
種別 & オープンウェイト \\
時系列 & 2024-11-26 \\
\end{tabularx}
\smallskip
{\bfseries 一次資料から整理したtechnical points}
\begin{itemize}[leftmargin=1.5em,itemsep=0.35em]
\item \textbf{一次情報で確認できる事実}: Ai2の一次資料により、OLMo 2について2024-11-26にモデル公開を確認した。 対象event近傍の一次資料から 7B parameter scale / 13B parameter scale / 5T training tokens を確認できる。
\end{itemize}
\begin{minipage}{\linewidth}
% reader-facing Technical Notes generic boundary/source tail group
\begin{samepage}
{\bfseries 一次資料}
\begingroup\sloppy
\Urlmuskip=0mu plus 2mu\relax
\begin{itemize}[leftmargin=1.5em,itemsep=0.25em]
\item {\scriptsize\url{https://allenai.org/blog/olmo2}}
\end{itemize}
\endgroup
\end{samepage}
\end{minipage}
\endgroup
\end{technicalnote}
\medskip

\begin{technicalnote}{Tülu 3}{主要資料}
\begingroup
% reader-facing Technical Notes break policy
\widowpenalty=10000
\clubpenalty=10000
\displaywidowpenalty=10000
\begin{tabularx}{\linewidth}{@{}>{\bfseries}p{0.22\linewidth}X@{}}
組織 & Ai2 \\
種別 & Framework \\
時系列 & 2024-11-21 \\
\end{tabularx}
\smallskip
{\bfseries 一次資料から整理したtechnical points}
\begin{itemize}[leftmargin=1.5em,itemsep=0.35em]
\item \textbf{一次情報で確認できる事実}: Ai2の一次資料により、Tülu 3について2024-11-21にSoftware公開を確認した。 対象event近傍の一次資料から 70B parameter scale / verifiable rewards / open weights を確認できる。
\end{itemize}
\begin{minipage}{\linewidth}
% reader-facing Technical Notes generic boundary/source tail group
\begin{samepage}
{\bfseries 一次資料}
\begingroup\sloppy
\Urlmuskip=0mu plus 2mu\relax
\begin{itemize}[leftmargin=1.5em,itemsep=0.25em]
\item {\scriptsize\url{https://allenai.org/blog/tulu-3-technical}}
\end{itemize}
\endgroup
\end{samepage}
\end{minipage}
\endgroup
\end{technicalnote}
\medskip

\begin{technicalnote}{Jamba 1.5 Mini and Large}{主要資料}
\begingroup
% reader-facing Technical Notes break policy
\widowpenalty=10000
\clubpenalty=10000
\displaywidowpenalty=10000
\begin{tabularx}{\linewidth}{@{}>{\bfseries}p{0.22\linewidth}X@{}}
組織 & AI21 \\
種別 & オープンウェイト \\
時系列 & 2024-08-22 \\
\end{tabularx}
\smallskip
{\bfseries 一次資料から整理したtechnical points}
\begin{itemize}[leftmargin=1.5em,itemsep=0.35em]
\item \textbf{一次情報で確認できる事実}: AI21の一次資料により、Jamba 1.5 Mini and Largeについて2024-08-22にモデル公開を確認した。 AI21の受理済み2024-08-22一次資料では、Jamba 1.5 Mini / LargeをSSM-Transformer architecture上のopen model familyとして公開し、TransformerとMambaを組み合わせるhybrid構成、双方256K effective context、Jamba Open Model Licenseを明示する。developer-facing機能としてstructured JSON output、function calling、document objects、citationsを挙げる。Miniを単一GPUで最大140K tokensまで扱えるという記述はdeployment条件であり、model context windowそのもの（256K）とは区別する。Llama 3.1 70B / 405B等は比較対象であり、Jamba自身のparameter scaleへ昇格させない。
\end{itemize}
\begin{minipage}{\linewidth}
% reader-facing Technical Notes generic boundary/source tail group
\begin{samepage}
{\bfseries 一次資料}
\begingroup\sloppy
\Urlmuskip=0mu plus 2mu\relax
\begin{itemize}[leftmargin=1.5em,itemsep=0.25em]
\item {\scriptsize\url{https://www.ai21.com/blog/announcing-jamba-model-family/}}
\end{itemize}
\endgroup
\end{samepage}
\end{minipage}
\endgroup
\end{technicalnote}
\medskip

\begin{technicalnote}{Hunyuan-Large}{主要資料}
\begingroup
% reader-facing Technical Notes break policy
\widowpenalty=10000
\clubpenalty=10000
\displaywidowpenalty=10000
\begin{tabularx}{\linewidth}{@{}>{\bfseries}p{0.22\linewidth}X@{}}
組織 & Tencent Hunyuan \\
種別 & オープンウェイト \\
時系列 & 2024-11-05 \\
\end{tabularx}
\smallskip
{\bfseries 一次資料から整理したtechnical points}
\begin{itemize}[leftmargin=1.5em,itemsep=0.35em]
\item \textbf{一次情報で確認できる事実}: Tencent Hunyuanの一次資料により、Hunyuan-Largeについて2024-11-05にモデル公開を確認した。 Tencent Hunyuanの受理済み2024-11-05一次資料では、Hunyuan-Large（Hunyuan-MoE-A52B）をTransformer-based MoEとして公開し、389B total / 52B active parametersと記載している。pretrained modelは256K、Instruct modelは128Kのtext sequenceをサポートし、GQAとCross-Layer AttentionによるKV-cache圧縮、FP8版、vLLM backend、training scriptsとmodel implementationsを公開範囲として示す。throughput、memory削減、benchmark等の性能値はTencent Hunyuanの報告として扱い、比較対象モデルの数値をHunyuan-Large自身の仕様へ昇格させない。
\end{itemize}
\begin{minipage}{\linewidth}
% reader-facing Technical Notes generic boundary/source tail group
\begin{samepage}
{\bfseries 一次資料}
\begingroup\sloppy
\Urlmuskip=0mu plus 2mu\relax
\begin{itemize}[leftmargin=1.5em,itemsep=0.25em]
\item {\scriptsize\url{https://github.com/Tencent-Hunyuan/Tencent-Hunyuan-Large}}
\end{itemize}
\endgroup
\end{samepage}
\end{minipage}
\endgroup
\end{technicalnote}
\medskip

\begin{technicalnote}{Phi-3.5 model family}{主要資料}
\begingroup
% reader-facing Technical Notes break policy
\widowpenalty=10000
\clubpenalty=10000
\displaywidowpenalty=10000
\begin{tabularx}{\linewidth}{@{}>{\bfseries}p{0.22\linewidth}X@{}}
組織 & Microsoft \\
種別 & オープンウェイト \\
時系列 & 2024-08-22 \\
\end{tabularx}
\smallskip
{\bfseries 一次資料から整理したtechnical points}
\begin{itemize}[leftmargin=1.5em,itemsep=0.35em]
\item \textbf{一次情報で確認できる事実}: Microsoftの一次資料により、Phi-3.5 model familyについて2024-08-22にモデル公開を確認した。 Microsoftの受理済み一次資料では、Phi-3.5-mini（3.8B・128K context）、Phi-3.5-vision（multi-frame image understanding）、Phi-3.5-MoE（16 experts、42B total / 6.6B active・128K context）を同じfamilyの別variantとして説明している。これらの数値・能力表現はMicrosoft記載値として扱う。
\end{itemize}
\begin{minipage}{\linewidth}
% reader-facing Technical Notes generic boundary/source tail group
\begin{samepage}
{\bfseries 一次資料}
\begingroup\sloppy
\Urlmuskip=0mu plus 2mu\relax
\begin{itemize}[leftmargin=1.5em,itemsep=0.25em]
\item {\scriptsize\url{https://techcommunity.microsoft.com/blog/azure-ai-foundry-blog/discover-the-new-multi-lingual-high-quality-phi-3-5-slms/4225280/}}
\end{itemize}
\endgroup
\end{samepage}
\end{minipage}
\endgroup
\end{technicalnote}
\medskip

\begin{technicalnote}{Ministral 3B and 8B}{主要資料}
\begingroup
% reader-facing Technical Notes break policy
\widowpenalty=10000
\clubpenalty=10000
\displaywidowpenalty=10000
\begin{tabularx}{\linewidth}{@{}>{\bfseries}p{0.22\linewidth}X@{}}
組織 & Mistral AI \\
種別 & モデル \\
時系列 & 2024-10-16 \\
\end{tabularx}
\smallskip
{\bfseries 一次資料から整理したtechnical points}
\begin{itemize}[leftmargin=1.5em,itemsep=0.35em]
\item \textbf{一次情報で確認できる事実}: Mistral AIの一次資料により、Ministral 3B and 8Bについて2024-10-16にモデル公開を確認した。 Mistral AIの受理済み2024-10-16一次資料では、on-device / edge向けの対象modelをMinistral 3BとMinistral 8Bとして明示し、双方最大128K context（当時vLLMでは32K）をサポートすると記載する。Ministral 8Bはfaster / memory-efficient inference向けのspecial interleaved sliding-window attention patternを持つ。sub-10Bはカテゴリ表現であり、Gemma 2 2B、Mistral 7B、Gemma 2 9B等はbenchmark比較対象なので、対象model自身のparameter scaleへ昇格させない。
\end{itemize}
\begin{minipage}{\linewidth}
% reader-facing Technical Notes generic boundary/source tail group
\begin{samepage}
{\bfseries 一次資料}
\begingroup\sloppy
\Urlmuskip=0mu plus 2mu\relax
\begin{itemize}[leftmargin=1.5em,itemsep=0.25em]
\item {\scriptsize\url{https://mistral.ai/news/ministraux/}}
\end{itemize}
\endgroup
\end{samepage}
\end{minipage}
\endgroup
\end{technicalnote}
\medskip

\begin{technicalnote}{Mistral Large 2}{補足資料}
\begingroup
% reader-facing Technical Notes break policy
\widowpenalty=10000
\clubpenalty=10000
\displaywidowpenalty=10000
\begin{tabularx}{\linewidth}{@{}>{\bfseries}p{0.22\linewidth}X@{}}
組織 & Mistral AI \\
種別 & モデル \\
時系列 & 2024-07-24 \\
\end{tabularx}
\smallskip
{\bfseries 一次資料から整理したtechnical points}
\begin{itemize}[leftmargin=1.5em,itemsep=0.35em]
\item \textbf{一次情報で確認できる事実}: Mistral AIの一次資料により、Mistral Large 2について2024-07-24にモデル公開を確認した。 Mistral AIの受理済み2024-07-24一次資料では、Mistral Large 2自身を123B parameters・128K contextのmodelとして記載し、single-node inferenceを想定する。Mistral Research Licenseでresearch / non-commercial use向けweightsを提供し、la Plateformeではmistral-large-2407として利用できる。同modelの機能としてparallel / sequential function callingとretrieval skillsを説明する。Codestral 22B、Codestral Mamba、Llama 3 405Bは同ページ内の関連製品・比較対象であり、Mistral Large 2自身のparameter scaleやarchitectureへ昇格させない。
\end{itemize}
\begin{minipage}{\linewidth}
% reader-facing Technical Notes generic boundary/source tail group
\begin{samepage}
{\bfseries 一次資料}
\begingroup\sloppy
\Urlmuskip=0mu plus 2mu\relax
\begin{itemize}[leftmargin=1.5em,itemsep=0.25em]
\item {\scriptsize\url{https://mistral.ai/news/mistral-large-2407/}}
\end{itemize}
\endgroup
\end{samepage}
\end{minipage}
\endgroup
\end{technicalnote}
\medskip

\begin{technicalnote}{Qwen2.5-Coder family}{補足資料}
\begingroup
% reader-facing Technical Notes break policy
\widowpenalty=10000
\clubpenalty=10000
\displaywidowpenalty=10000
\begin{tabularx}{\linewidth}{@{}>{\bfseries}p{0.22\linewidth}X@{}}
組織 & Qwen \\
種別 & オープンウェイト \\
時系列 & 2024-11-12 \\
\end{tabularx}
\smallskip
{\bfseries 一次資料から整理したtechnical points}
\begin{itemize}[leftmargin=1.5em,itemsep=0.35em]
\item \textbf{一次情報で確認できる事実}: Qwenの一次資料により、Qwen2.5-Coder familyについて2024-11-12にモデル公開を確認した。 Qwen Teamの受理済み2024-11-12一次資料のmodel tableでは、Qwen2.5-Coderを0.5B / 1.5B / 3B / 7B / 14B / 32Bのvariantとして列挙している。licenseは0.5B / 1.5B / 7B / 14B / 32BがApache 2.0、3BがQwen Researchであり、family-levelの単一licenseとして平坦化しない。
\end{itemize}
\begin{minipage}{\linewidth}
% reader-facing Technical Notes generic boundary/source tail group
\begin{samepage}
{\bfseries 一次資料}
\begingroup\sloppy
\Urlmuskip=0mu plus 2mu\relax
\begin{itemize}[leftmargin=1.5em,itemsep=0.25em]
\item {\scriptsize\url{https://qwenlm.github.io/blog/qwen2.5-coder-family/}}
\end{itemize}
\endgroup
\end{samepage}
\end{minipage}
\endgroup
\end{technicalnote}
\medskip

\begin{technicalnote}{Qwen2.5-Turbo 1M context}{補足資料}
\begingroup
% reader-facing Technical Notes break policy
\widowpenalty=10000
\clubpenalty=10000
\displaywidowpenalty=10000
\begin{tabularx}{\linewidth}{@{}>{\bfseries}p{0.22\linewidth}X@{}}
組織 & Qwen \\
種別 & モデル更新 \\
時系列 & 2024-11-15 \\
\end{tabularx}
\smallskip
{\bfseries 一次資料から整理したtechnical points}
\begin{itemize}[leftmargin=1.5em,itemsep=0.35em]
\item \textbf{一次情報で確認できる事実}: Qwenの一次資料により、Qwen2.5-Turbo 1M contextについて2024-11-15にモデル公開を確認した。 Qwen Teamの受理済み2024-11-15一次資料では、Qwen2.5-Turboのcontext lengthを128Kから1Mへ拡張したと記載し、extremely long context向けにsparse attentionを用いた推論最適化を説明している。Alibaba Cloud Model StudioのAPI serviceに加えてHugging Face Demo / ModelScope Demoで提供された。1M contextでのTTFT 4.9分→68秒（4.3x）、Passkey/RULER等の性能値はQwen Team自身の報告として扱い、独立再現結果へ昇格させない。
\end{itemize}
\begin{minipage}{\linewidth}
% reader-facing Technical Notes generic boundary/source tail group
\begin{samepage}
{\bfseries 一次資料}
\begingroup\sloppy
\Urlmuskip=0mu plus 2mu\relax
\begin{itemize}[leftmargin=1.5em,itemsep=0.25em]
\item {\scriptsize\url{https://qwenlm.github.io/blog/qwen2.5-turbo/}}
\end{itemize}
\endgroup
\end{samepage}
\end{minipage}
\endgroup
\end{technicalnote}
\medskip

\begin{technicalnote}{Quantized Llama 3.2 lightweight models}{補足資料}
\begingroup
% reader-facing Technical Notes break policy
\widowpenalty=10000
\clubpenalty=10000
\displaywidowpenalty=10000
\begin{tabularx}{\linewidth}{@{}>{\bfseries}p{0.22\linewidth}X@{}}
組織 & Meta \\
種別 & モデル更新 \\
時系列 & 2024-10-24 \\
\end{tabularx}
\smallskip
{\bfseries 一次資料から整理したtechnical points}
\begin{itemize}[leftmargin=1.5em,itemsep=0.35em]
\item \textbf{一次情報で確認できる事実}: Metaの一次資料により、Quantized Llama 3.2 lightweight modelsについて2024-10-24にモデル公開を確認した。 Metaの受理済み2024-10-24一次資料は、Llama 3.2 1B/3B instruction-tuned modelsの量子化版をmobile / edge向けに公開した。QAT + LoRA（記事中ではQLoRAと呼称）とSpinQuantを用い、linear layersは4-bit groupwise weights + 8-bit per-token dynamic activations、classification layerは8-bit per-channel weights + 8-bit per-token dynamic activations、embeddingは8-bit per-channel quantizationと記載している。ExecuTorch / Arm CPUを主対象とし、short-context applicationsは最大8Kを優先する。2–4x speedup、model size平均56\%減、memory平均41\%減等はMetaのOnePlus 12等での測定値として扱い、独立再現結果へ昇格させない。
\end{itemize}
\begin{minipage}{\linewidth}
% reader-facing Technical Notes generic boundary/source tail group
\begin{samepage}
{\bfseries 一次資料}
\begingroup\sloppy
\Urlmuskip=0mu plus 2mu\relax
\begin{itemize}[leftmargin=1.5em,itemsep=0.25em]
\item {\scriptsize\url{https://ai.meta.com/blog/meta-llama-quantized-lightweight-models/}}
\end{itemize}
\endgroup
\end{samepage}
\end{minipage}
\endgroup
\end{technicalnote}
\medskip
